\id{IRSTI 28.23.15}{}

\begin{header}
\swa{T.K.~Azbergen, A.O.~Tleubayeva, A.~Aituarov, Y.T.~Sabyrzhan}{BENCHMARKING CALIBRATED GRADIENT BOOSTING WITH BEHAVIOURAL FEATURES FOR TRANSACTION-LEVEL AML DETECTION}

\tsp{1}T.K.~Azbergen,
\tsp{2}A.O.~Tleubayeva\envelope,
\tsp{1}A.T~Aituarov,
\tsp{1}Y.T.~Sabyrzhan
\end{header}

\begin{affil}
\tsp{1}Kazakh-British Technical University, Almaty, Kazakhstan,

\tsp{2}Astana IT University, Astana, Kazakhstan

\corrauthor{Corresponding-author: Arailym.tll@gmail.com}
\end{affil}

Reliable transaction-level anti-money-laundering (AML) systems require
not only accurate ranking of suspicious activity but also
well-calibrated probabilistic risk estimates to support operational
decision-making and risk prioritisation. In this context, this work
provides an empirical evaluation of calibrated gradient boosting models
with temporal behavioural features, focusing on the combined impact of
probabi\-listic calibration and behavioural feature modelling on
probability reliability and operationally relevant performance metrics.

The proposed framework integrates gradient boosting algorithms (CatBoost
and XGBoost) with post-hoc probability calibration techniques (Platt
scaling and isotonic regression) and behavioural features computed over
rolling time windows of 1, 7, and 30 days to capture both short-term
volatility and longer-term transaction patterns. Experimental evaluation
is conducted on the publicly available IBM Anti-Money Laundering
(IBM--AML) dataset, which contains approximately seven million synthetic
transac\-tions simulating realistic banking activity.

Model performance is assessed using ranking-, calibration-, and
operationally oriented metrics, includ\-ing AUPRC, AUROC, Brier score,
Precision@Top k, and Lift@Top k. Among the evaluated baselines, the
calibrated CatBoost model showed the strongest overall performance among
the evaluated baselines (AUPRC = 0.367, AUROC = 0.959, Brier = 0.0062,
Lift@1\% = 40.9).

The results indicate that probability calibration improves the
reliability of predicted risk scores without degrading ranking
performance, while temporal behavioural features contribute to improved
detection sensitivity in highly imbalanced settings. This study provides
empirical evidence and methodological clarity regarding the role of
calibration and behavioural dynamics in AML systems and offers a
benchmark-ready кexperimental pipeline for future research. The
applicability of the findings to real-world AML systems is discussed in
light of the synthetic nature of the dataset, and directions for further
validation on real banking data are outlined.

{\bfseries Keywords:}~machine learning, anti-money laundering, CatBoost,
probabilistic calibration, behavioural features, financial fraud
detection, explainable AI.

\begin{header}
СРАВНИТЕЛЬНАЯ ОЦЕНКА КАЛИБРОВАННОГО ГРАДИЕНТНОГО БУСТИНГА С ПОВЕДЕНЧЕСКИМИ ПРИЗНАКАМИ ДЛЯ ВЫЯВЛЕНИЯ ОТМЫВАНИЯ ДЕНЕГ НА УРОВНЕ ТРАНЗАКЦИЙ

\tsp{1}Т.К.~Азберген,
\tsp{2}А.О.~Тлеубаева\envelope,
\tsp{1}А.Т.~Айтуов,
\tsp{1}Е.Т.~Сабыржан
\end{header}

\begin{affil}
\tsp{1}Казахстанско-Британский технический университет, Алматы, Казахстан

\tsp{2}Astana IT University, Астана, Казахстан

e-mail: Arailym.tll@gmail.com
\end{affil}

Надёжные системы противодействия отмыванию денег (AML) на уровне
отдельных транзакций требуют не только точного ранжирования
подозрительной активности, но и хорошо откалиброванных вероятностных
оценок риска для поддержки операционных решений и приоритизации
проверок. В этом контексте в работе представлена эмпирическая оценка
калиброванных моделей градиентного бустинга с временными поведенческими
признаками, направленная на анализ совместного влияния вероятностной
калибровки и моделирования поведенческих признаков на надёжность
вероятностных оценок и операционно значимые показатели эффективности.

Предложенная структура объединяет алгоритмы градиентного бустинга
(CatBoost и XGBoost) с методами пост-калибровки вероятностей (Platt
scaling и изотоническая регрессия), а также поведенческие признаки,
вычисляемые в скользящих временных окнах продолжительностью 1, 7 и 30
дней, что позволяет учитывать, как краткосрочную волатильность, так и
долгосрочные транзакционные паттерны. Экспериментальная оценка проведена
на открытом наборе данных IBM Anti-Money Laundering (IBM--AML),
содержащем около семи миллионов синтетических транзакций, имитирующих
реалистичную банковскую деятельность.

Качество моделей оценивается с использованием ранжирующих, калибровочных
и операционно ориентированных метрик, включая AUPRC, AUROC, Brier score,
Precision@Top k и Lift@Top k. Среди рассмотренных базовых моделей
калиброванная модель CatBoost продемонстрировал самые высокие общие
показатели среди оцененных базовых показателей (AUPRC = 0.367, AUROC =
0.959, Brier = 0.0062, Lift@1\% = 40.9).

Полученные результаты показывают, что вероятностная калибровка повышает
надёжность прогнозируемых риск-оценок без ухудшения качества
ранжирования, в то время как временные поведенческие признаки
способствуют повышению чувствительности выявления в условиях сильного
дисбаланса классов. Работа предоставляет эмпирические доказательства и
методологическую ясность относительно роли калибровки и поведенческой
динамики в AML-системах, а также предлагает воспроизводимую
экспериментальную основу для дальнейших исследований. Применимость
полученных выводов к реальным AML-системам обсуждается с учётом
синтетической природы используемого набора данных, а также обозначаются
направления дальнейшей валидации на реальных банковских данных.

{\bfseries Ключевые слова:} машинное обучение, отмывание денег, CatBoost,
вероятностная калибровка, поведенческие признаки, финансовое
мошенничество, объяснимая ИИ.

\begin{header}
ТРАНЗАКЦИЯЛАР ДЕҢГЕЙІНДЕ АҚШАНЫ ЖЫЛЫСТАТУДЫ АНЫҚТАУҒА АРНАЛҒАН МІНЕЗ-ҚҰЛЫҚТЫҚ БЕЛГІЛЕРГЕ НЕГІЗДЕЛГЕН КАЛИБРЛЕНГЕН ГРАДИЕНТТІ БУСТИНГТІҢ САЛЫСТЫРМАЛЫ БАҒАЛАНУЫ

\tsp{1}Т.К.~Азберген,
\tsp{2}А.О.~Тлеубаева\envelope,
\tsp{1}А.Т.~Айтуов,
\tsp{1}Е.Т.~Сабыржан
\end{header}

\begin{affil}
\tsp{1}Қазақ-Британ техникалық университеті, Алматы, Қазақстан,
       
\tsp{2}Astana IT University, Астана, Казахстан,

e-mail: Arailym.tll@gmail.com
\end{affil}

Транзакциялар деңгейіндегі ақшаны жылыстатуға қарсы (AML) сенімді
жүйелер күдікті әрекеттерді дәл ранжирлеумен қатар, операциялық шешім
қабылдау мен тәуекелдерді басымдықпен қарау үшін жақсы калибрленген
ықтималдықтық тәуекел бағаларын талап етеді. Осы тұрғыда бұл жұмыста
уақыттық мінез-құлықтық белгілерді қолданатын калибрленген градиентті
бустинг модельдерінің эмпирикалық бағалануы ұсынылады және ықтималдық
калибрлеу мен мінез-құлықтық белгілерді модельдеудің ықтималдық
бағалардың сенімділігіне және операциялық тұрғыдан маңызды тиімділік
көрсеткіштеріне бірлескен әсері талданады.

Ұсынылған құрылым градиентті бустинг алгоритмдерін (CatBoost және
XGBoost) ықтималдықтарды пост-калибрлеу әдістерімен (Platt scaling және
изотондық регрессия) біріктіреді, сондай-ақ 1, 7 және 30 күндік
жылжымалы уақыттық терезелерде есептелетін мінез-құлықтық белгілерді
пайдаланады. Бұл тәсіл қысқа мерзімді құбылмалылықты да, ұзақ мерзімді
транзакциялық үлгілерді де қамтуға мүмкіндік береді. Эксперименттік
бағалау нақты банктік қызметті модельдейтін шамамен жеті миллион
синтетикалық транзакциядан тұратын IBM Anti-Money Laundering (IBM--AML)
ашық деректер жиынтығында жүргізілді.

Модельдердің тиімділігі ранжирлеу, калибрлеу және операциялық
бағытталған метрикалар арқылы бағаланды, олардың қатарына AUPRC, AUROC,
Brier көрсеткіші, Precision@Top k және \\Lift@Top k жатады. Қарастырылған
базалық модельдер арасында калибрленген CatBoost бағалан\-ған базалық
көрсеткіштер арасында ең жоғары жалпы көрсеткіштерді көрсетті (AUPRC =
0.367, AUROC = 0.959, Brier = 0.0062, Lift@1\% = 40.9).

Нәтижелер ықтималдық калибрлеудің модельдің ранжирлеу қабілетін
төмендетпей, болжанған тәуекел бағаларының сенімділігін арттыратынын, ал
уақыттық мінез-құлықтық белгілердің сыныптар қатты теңгерімсіз болған
жағдайда анықтау сезімталдығын күшейтетінін көрсетеді. Зерттеу AML
жүйелеріндегі калибрлеу мен мінез-құлықтық динамиканың рөліне қатысты
эмпирикалық дәлелдер мен әдістемелік айқындықты ұсынады және болашақ
зерттеулер үшін қайта өндіруге болатын эксперименттік негіз
қалыптастырады. Алынған нәтижелердің нақты AML жүйелеріне қолданылу
мүмкіндігі пайдаланылған деректер жиынтығының синтетикалық сипатын
ескере отырып талқыланып, нақты банктік деректерде қосымша валидация
жүргізу бағыттары айқындалады.

{\bfseries Түйін сөздер:~}машиналық оқыту, ақшаны жылыстату, CatBoost,
ықтималдық калибрлеу, мінез-құлықтық ерекшеліктер, қаржылық алаяқтық,
түсіндірілетін жасанды интеллект.

\begin{multicols}{2}
{\bfseries Introduction.} Money laundering (AML) remains one of the most
persistent challenges to global financial stability, affecting both the
integrity of banking institutions and the credibility of regulatory
systems. According to the~Financial Action Task Force (FATF),
illicit financial flows represent approximately~2-5\% of global
GDP, illustrating the magnitude of the threat and the necessity
for advanced analytical tools to counteract it. The rapid evolution of
laundering schemes, the expansion of digital banking platforms, and the
increasing complexity of transactional ecosystems have made the
detection of suspicious activity a pressing global priority {[}1{]}.

Traditional~rule-based AML systems, which have long dominated industry
practice, rely on static thresholds, expert-defined rules, and manually
specified behavioural heuristics {[}2{]}. While these systems are
transparent and straightforward to audit, they exhibit several critical
limitations. They tend to produce excessively high false-positive
rates-often exceeding~90\%~of all alerts {[}3{]}, which leads to a heavy
compliance burden, inflated operational costs, and reduced confidence in
automated systems. Moreover, rule-based models struggle to adapt to new
or evolving laundering typologies, as their decision logic remains fixed
once deployed.

In contrast,~machine learning (ML)~offers a data-driven and adaptive
alternative capable of identifying complex behavioural and temporal
patterns in financial transactions {[}4{]}. ML-based approaches enable
dynamic adjustment to new data distributions, capture non-linear
relationships, and provide improved detection performance across
heterogeneous client profiles and transaction types {[}5{]}. This shift
reflects a broader paradigm in financial monitoring: from deterministic
rules toward statistically grounded, probabilistically interpretable
models that learn from empirical evidence.

Recent studies have highlighted the effectiveness of supervised ML
models in AML detection, frequently reporting that gradient boosting
methods-such as XGBoost, LightGBM, and CatBoost-outperform traditional
rule-based heuristics in identifying suspicious transactions. However,
reported performance gains are often highly sensitive to feature
construction strategies, data-splitting protocols, and the choice of
evaluation metrics. In particular, much of the existing literature
prioritises global discrimination measures, such as the area under the
receiver operating characteristic curve (ROC-AUC) and the F1-score,
while giving comparatively limited attention to probability reliability
and operational prioritization {[}6{]}.

As a result, despite achieving strong ranking performance, many AML
models produce risk scores whose probabilistic interpretation is unclear
or unreliable, limiting their usefulness for risk-based decision-making
in compliance workflows. Two aspects are especially critical in this
regard. First, probability calibration ensures that predicted risk
scores correspond to empirically valid probabilities, enabling
consistent prioritisation of alerts under constrained review resources.
Second, temporal behavioural representation captures the dynamic
evolution of client activity over time, rather than relying on static
transaction snapshots, thereby improving sensitivity to subtle and
emerging laundering patterns.

In parallel, graph-based learning approaches have been explored to
represent financial systems as networks of interconnected accounts,
entities, and transactions. Graph Neural Networks (GNNs) {[}7{]},
including Graph Convolutional Networks (GCNs) {[}8{]}, dynamic GNN
architectures such as MDGC-LSTM {[}9{]}, and multi-view formulations
such as LineMVGNN {[}10{]}, have demonstrated promise in modelling
multi-stage laundering schemes and uncovering hidden relational
dependencies. Nevertheless, many graph-based AML models rely on
uncalibrated outputs and provide limited temporal context, which can
constrain their interpretability, stability, and suitability for
operational compliance settings.

Compared to prior work, existing behavioural ML approaches {[}6{]}
primarily emphasise feature design and anomaly sensitivity, but
typically evaluate models using discrimination-focused metrics without
explicitly addressing probability calibration or operational
prioritisation. At the same time, recent graph-based AML studies
{[}7-10{]} demonstrate the value of relational representations for
uncovering complex laundering structures, yet often rely on uncalibrated
outputs and lack systematic temporal aggregation aligned with real
monitoring workflows. As a result, probability reliability and risk
score interpretability-both critical for compliance
decision-making-remain underexplored in much of the existing literature.

Taken together, these observations highlight a clear research gap.
Existing AML models rarely combine temporal behavioural modelling,
probability calibration, and graph-based contextual learning within a
unified, reproducible framework. Consequently, this study aims to
empirically investigate this gap by evaluating how calibration-aware and
behaviour-aware machine-learning pipelines perform under operationally
relevant AML metrics, rather than by proposing a new detection paradigm.

Accordingly, this work is positioned as an empirical evaluation study.
The main contribution is an empirically grounded evaluation protocol
that jointly analyses calibration, behavioural temporal features, and
operational top-k metrics within a reproducible benchmark setting.

The object of this research is the process of detecting and classifying
suspicious financial transactions within an AML system. The subject of
the research is the application of calibrated and behaviour-aware
machine-learning methods for improving the reliability of AML risk
scoring. The purpose of the study is to empirically assess an integrated
AML framework that combines temporal behavioural features and
probabilistic calibration under ranking- and operation-oriented
evaluation criteria.

This work is based on the hypothesis that the integration of temporal
behavioural features and probability calibration can improve ranking
stability and probability reliability in highly imbalanced AML detection
tasks. Specifically, calibrated boosting algorithms such as CatBoost,
when combined with rolling behavioural features, can yield more stable
and better-calibrated AML risk scores than uncalibrated or purely
graph-based models.

Accordingly, the study pursues the following objectives:

To critically analyse the limitations of rule-based and conventional ML
approaches in AML detection.

- To design a~behavioural feature-engineering pipeline~that captures
temporal transaction dynamics through rolling windows of 1, 7, and 30
days.

- To implement and compare~probability calibration techniques,
including Platt scaling and isotonic regression.

- To evaluate and contrast the performance of~calibrated boosting models
(CatBoost, XGBoost)~and~graph-based architectures
(EdgeSAGE)~on a unified dataset.

- To develop a~reproducible experimental framework~based on the~IBM-AML
dataset, enabling transparent benchmarking for future
research.

The methodological framework combines supervised
learning, probability calibration, and~graph-based analysis.
Feature engineering incorporates rolling behavioural statistics to
capture temporal variability in transaction activity.

The dataset employed - the~IBM Anti-Money Laundering (AML)~dataset
provides a realistic synthetic simulation of financial transactions
{[}11{]}. Model performance is assessed using~ranking metrics~(AUPRC,
AUROC) {[}12{]},~calibration metrics (Brier score, reliability curves)
{[}13{]}, and~operational metrics~(Precision@Top k, Lift@Top k)
{[}14{]}. All experiments follow a~temporal data-splitting strategy~to
preserve chronological order and simulate realistic monitoring
scenarios.

The scientific novelty of this study~does not consist in the development
of a fundamentally new machine-learning algorithm. Instead, it lies in a
structured and reproducible empirical analysis of how temporal
behavioural modelling and probabilistic calibration jointly affect
reliability and operational interpretability in AML detection tasks.

Unlike prior research focused primarily on discrimination accuracy, this
work emphasises probability reliability and operational evaluation
metrics as complementary dimensions of AML model assessment. This focus
is motivated by the practical need to support risk prioritisation and
auditability, where relative ranking alone is insufficient without
reliable probability estimates.

The theoretical significance of the study lies in~advancing empirical
understanding of calibration-aware and behaviour-aware machine-learning
approaches in financial compliance.

The practical significance is reflected in a benchmark-ready evaluation
pipeline that can support future comparative studies and validation on
real-world financial data.

{\bfseries Materials and methods.} The experiments were conducted using
the~IBM Anti-Money Laundering (AML) dataset, a publicly available
synthetic dataset designed to emulate real-world financial transactions.
It contains~approximately 7 million transaction records, representing a
virtual banking environment with both legitimate and illicit financial
activities. Each transaction includes structured attributes describing
the sender, receiver, transaction amount, currency, payment type, and
temporal information.

Since the original dataset is highly imbalanced, with an approximate
ratio of 1:200 between suspicious and non-suspicious transactions, a
stratified sampling procedure was applied prior to model training to
reduce the dataset size while preserving the original class
distribution. Temporal ordering was maintained during the
train-validation-test split to prevent information leakage, resulting in
a manageable and statistically representative subset of approximately
1.5 million records suitable for training and evaluation.

To maintain temporal consistency and reflect real monitoring conditions,
a~chronological data split~was employed:

- 70\%~of transactions for training,

- 15\%~for validation, and

- 15\%~for testing.

This temporal partition prevents information leakage and allows for
backtesting across monthly periods (e.g., 2022-09) to assess model
stability over time.

Feature engineering is central to this study and aims to
capture~behavioural and temporal patterns~in transaction dynamics.

Three major feature categories were developed:

1. Categorical features~- nominal attributes describing transaction type,
currency, bank identifier, and payment format (≈5 variables). These
were encoded using target or one-hot encoding depending on
cardinality.

2. Quantitative features~- numerical indicators such as transaction
amount, average, standard deviation, and count of transactions.

3. Behavioural rolling features - temporal aggregates computed over 1-,
7-, and 30-day sliding windows, including:

- mean and median transaction values;

- transaction count per window;

- rolling sum and deviation (MAD, standard deviation) {[}15{]}.

These features represent~short-term and long-term behavioural trends,
allowing the model to detect deviations from normal patterns that may
indicate laundering activities.\\
All numerical features were scaled using~z-score normalization,
and missing values were handled through median imputation.

The proposed AML detection framework integrates both~classical
machine-learning models~and~graph-based architectures, enabling
comparative evaluation of their predictive and operational
effectiveness.

Four supervised learning algorithms were trained on the tabular feature
set:

1. Logistic Regression (LR)~- interpretable linear baseline;

2. Random Forest (RF)~- ensemble of decision trees for non-linear
modelling;

3. XGBoost~- gradient boosting framework with efficient regularization;

4. CatBoost~- categorical boosting model optimized for mixed data types.

CatBoost and XGBoost were selected as primary baselines because gradient
boosting is a widely adopted and highly effective model class for
heterogeneous tabular data, including imbalanced classification settings
{[}16, 17{]}. CatBoost additionally provides built-in mechanisms for
handling categorical features and was designed to reduce prediction
shift caused by target leakage during boosting, which is beneficial in
mixed-type transaction data {[}18{]}.

We focused on CatBoost and XGBoost as representative, widely deployed
gradient-boosting implementations with mature open-source ecosystems and
strong empirical performance on heterogeneous tabular data. LightGBM was
not included to keep the benchmark compact and avoid overlapping
baselines with similar inductive bias under the same feature set; its
inclusion is left for future extensions of the benchmark. Neural tabular
models were deprioritised due to higher optimisation complexity and
stronger dependence on tuning choices, which can increase training
variance in operational pipelines.

Alternative deep neural models were not prioritised due to higher
optimisation complexity and stronger dependence on training and tuning
choices, which can increase variance and computational cost in practical
tabular pipelines {[}19, 20{]}. Graph neural networks were included as a
conceptual baseline for relational learning over account--transaction
graphs, following standard inductive graph representation learning
frameworks {[}21, 22{]}.

Each model was trained using a binary classification objective
(isMoneyLaundering = 1 vs.0), with class weights adjusted to counter
data imbalance.~

Hyperparameters were optimized on the validation set via grid search,
with early stopping based on validation AUPRC. To reduce sensitivity to
hyperparameter choice, all models were evaluated using fixed
validation-based configurations, and performance stability was assessed
through temporal backtesting and feature ablation experiments.~

To assess robustness with respect to hyperparameter choice, the
following tuning policy was adopted. Hyperparameter search was
intentionally constrained to a compact grid to reduce variance and avoid
overfitting to a single validation period, in line with established
recommendations for ensemble models {[}23{]}. Gradient boosting methods
are known to exhibit relatively low sensitivity to moderate variations
in depth and learning rate due to their ensemble structure {[}17{]}. In
our experiments, CatBoost and XGBoost demonstrated stable performance
across reasonable hyperparameter ranges, with only minor fluctuations in
AUPRC and top-k operational metrics, while post-hoc probability
calibration consistently improved the Brier score, as expected from
prior calibration studies {[}17,24{]}.

To enhance interpretability and improve risk-based ranking, model
outputs were calibrated using two standard post-processing techniques:

1. Platt Scaling (Sigmoid Calibration) - fits a logistic
regression model on validation predictions {[}25{]};

2. Isotonic Regression~- non-parametric monotonic mapping to align
predicted probabilities with empirical frequencies {[}24{]}.

3. Calibration models were fitted on the validation set and applied to
test predictions to ensure unbiased evaluation.

Calibration was applied to improve probability reliability and was
evaluated using the Brier score and reliability diagrams in the Results
section.

In addition to tabular ML models, a graph-based~EdgeSAGE~model
was implemented. The transaction network was constructed by connecting
nodes representing accounts through edges representing transactions
{[}22{]}.

Each edge carried feature vectors derived from the transaction
attributes. The EdgeSAGE model aggregated information from neighbouring
nodes to predict the likelihood of a transaction being suspicious.

The EdgeSAGE model was included as a graph-based baseline to enable
comparative evaluation with tabular models under the same experimental
conditions.

Model performance was evaluated using three complementary groups of
metrics designed to capture ranking quality, probabilistic calibration,
and operational effectiveness.

Given the strong class imbalance of the AML dataset (approximately
1:200), ranking and precision-based metrics were prioritised over
accuracy-based measures to provide a more realistic assessment of
detection capability.

Model calibration was performed using Platt scaling {[}25{]} and
isotonic regression {[}24{]}, fitted on validation predictions.

The~AUPRC~quantifies the trade-off between~precision~and~recall~across
varying decision thresholds, making it particularly suitable for highly
imbalanced datasets such as AML detection.

Precision (P) and Recall (R) are defined as:

\begin{equation}
\text{Precision} = \frac{TP}{TP + FP}, \quad \text{Recall} = \frac{TP}{TP + FN}
\end{equation}

where:

- \(TP\) - number of true positives (correctly identified suspicious
transactions),

- \(FP\) - number of false positives,

- \(FN\) - number of false negatives.

The area under the curve (AUPRC) is then computed as the integral of
precision with respect to recall:

\begin{equation}
\text{AUPRC} = \int_{0}^{1}{P(R)}\, dR
\end{equation}

Higher AUROC indicates better ranking of positive (suspicious) instances
among top risk scores.

The~AUROC~evaluates the model's ability to discriminate between positive
and negative classes across all thresholds. It is defined as the
probability that a randomly chosen positive instance receives a higher
predicted score than a randomly chosen negative instance:

\begin{equation}
\text{AUROC} = \int_{0}^{1}{TPR(FPR)}\, dFPR
\end{equation}

where:

\begin{equation}
TPR = \frac{TP}{TP + FN},\quad FPR = \frac{FP}{FP + TN}
\end{equation}

and \(TN\) denotes true negatives. While AUROC is a standard indicator
of model discrimination, AUPRC is generally more informative for
imbalanced AML data where false negatives are critical.

The~Brier score (BS)~measures the mean squared difference
between predicted probabilities \((p_{i})\) and actual outcomes
\((y_{i})\):

\begin{equation}
\text{Brier score} = \frac{1}{N}i = 1N\left( \widehat{p_{i}} - y_{i} \right)^{2}
\end{equation}

where:

- \(p_{i} \in \lbrack 0,1\rbrack\) is the predicted probability of
laundering for transaction~\emph{i,}

- \(y_{i} \in \text{\{}0,1\text{\}}\) is the ground-truth label,

- \(N\) is the number of samples

A lower Brier score indicates better~probabilistic reliability~and
calibration of the model.

To visually assess calibration, a reliability (calibration)
curve plots the average predicted probability against the
observed frequency of positive cases in corresponding probability bins.

Perfect calibration corresponds to the diagonal y=x, y=x.

Deviations from this line indicate overconfidence (predicted
probabilities too high) or underconfidence (too low). In AML
contexts, well-calibrated risk scores are essential for prioritising
transactions according to realistic probability estimates rather than
arbitrary ranking.

To reflect resource-constrained review settings, models are evaluated
under top-k selection scenarios, where only the highest-ranked
transactions are examined.

Precision@Topk measures the proportion of true positives within the
top \emph{k\%} of transactions with the highest predicted scores:

\begin{equation}
\text{PrecisionTopk} = TP_{k}TPk + FPk
\end{equation}

where~\(TPk\)~and~\(FPk\)~denote true and false
positives among the top~\(k\%\)~ranked cases.

\emph{Recall@Top k} measures how many of the actual laundering cases are
captured in the top~\(k\%\)~predictions:

\begin{equation}
\text{RecallTopk} = TP_{k}TP
\end{equation}

This metric reflects the system's ability to~capture critical
alerts within the operational review budget (e.g., the top 1\% of
transactions reviewed daily by analysts).

The Lift metric quantifies how much better the model performs than
random selection at a given cutoff. It is computed as the ratio between
the model's Precision@Top k and the base rate of positives in the
dataset:
\end{multicols}

\begin{equation}
LiftatTopk = \frac{\text{Precision at Top}\, k}{P_{\text{base}}},
\quad P_{\text{base}} = \frac{TP + FP}{N}
\end{equation}

\begin{multicols}{2}
Higher Lift@k values indicate a greater concentration of true positive
cases within the top-ranked transactions compared to random selection.
In applied AML evaluation scenarios, Lift@1\% values above 30× are
commonly interpreted as indicating a high concentration of detected
cases within a limited review subset.

{\bfseries Results and discussion.} The experimental evaluation was
conducted on the IBM-AML dataset {[}11{]}, comprising several million
synthetic financial transactions that simulate realistic banking
behaviour, including both legitimate and laundering operations.

The performance of multiple machine-learning models -Logistic
Regression, Random Forest, XGBoost, CatBoost, and a Graph Neural Network
(EdgeSAGE) - was assessed using three complementary groups of metrics:
ranking, calibration, and operational effectiveness.

Table 1 summarises the comparative performance of all baseline models
evaluated on the IBM-AML dataset.

In contrast to prior studies focusing primarily on ROC-AUC or F1-score
optimisation {[}3,5{]}, the present results indicate that
calibration-aware models can improve probability reliability without
sacrificing ranking performance.

This observation complements behavioural feature-based approaches
{[}6{]} by quantifying calibration effects under operational AML
metrics.
\end{multicols}

\tcap{Table 1 - Performance comparison of baseline models on the IBM-AML dataset}
\begin{longtblr}[
  label = none,
  entry = none,
]{
  width = \linewidth,
  colspec = {Q[317]Q[85]Q[87]Q[75]Q[152]Q[123]Q[98]},
  cells = {c},
  cells = {font = \small},
  row{1} = {font=\bfseries\small},
  row{5} = {font=\bfseries\small},
  hlines,
  vlines,
}
Model                           & AUPRC  & AUROC & Brier  & Precision@1\% & Recall@1\% & Lift@1\% \\
Logistic Regression             & 0.0479 & 0.894 & 0.0067 & 0.069         & 0.099      & 9.9×     \\
Random Forest                   & 0.0904 & 0.913 & 0.0066 & 0.153         & 0.219      & 21.9×    \\
XGBoost (cal.)                  & 0.1680 & 0.946 & 0.0062 & 0.196         & 0.282      & 28.2×    \\
CatBoost (cal.)                 & 0.3674 & 0.959 & 0.0062 & 0.285         & 0.409      & 40.9×    \\
Graph Neural Network (EdgeSAGE) & 0.0110 & 0.666 & 0.772  & 0.007         & 0.009      & 0.9×     
\end{longtblr}

\begin{multicols}{2}
The CatBoost model with Platt calibration achieved the strongest overall
performance among the evaluated models under the selected ranking,
calibration, and operational metrics, with an AUPRC of
0.367 and Lift@1\% of 40.9×, indicating that a large proportion
of true laundering cases is concentrated within the top-ranked
transactions.

This corresponds to an approximately four-fold improvement in
operational precision relative to XGBoost and a substantial gain over
random selection under the chosen top-k operational setting. The Brier
score of 0.0062 indicates strong probabilistic calibration, reflecting
improved alignment between predicted risk scores and observed event
frequencies. These estimates are specific to the IBM-AML benchmark and
may not directly transfer to real-world deployments.

In contrast, the Graph Neural Network (EdgeSAGE) underperformed
substantially (AUPRC = 0.011, AUROC = 0.666), primarily due to the
limited depth of neighbourhood aggregation and absence of temporal
contextualisation.

This result highlights the critical role of behavioural and temporal
dynamics in AML detection, which conventional GNNs, when applied without
explicit temporal modelling, may not adequately capture. Table
2 presents a focused comparison of CatBoost models before and
after probability calibration
\end{multicols}

\tcap{Table 2 - Effect of probability calibration on CatBoost performance}
\begin{longtblr}[
  label = none,
  entry = none,
]{
  width = \linewidth,
  colspec = {Q[302]Q[102]Q[88]Q[187]Q[148]Q[113]},
  cells = {c},
  row{1} = {font=\bfseries\small},
  cell{3}{2} = {font=\bfseries\small},
  cell{3}{3} = {font=\bfseries\small},
  cell{3}{4} = {font=\bfseries\small},
  cell{3}{5} = {font=\bfseries\small},
  cell{3}{6} = {font=\bfseries\small},
  cell{4}{3} = {font=\bfseries\small},
  hlines,
  vlines,
}
Model                       & AUPRC  & Brier  & Precision@1\% & Recall@1\% & Lift@1\% \\
Uncalibrated CatBoost       & 0.3674 & 0.1031 & 0.285         & 0.409      & 40.9×    \\
Platt (Sigmoid) Calibration & 0.3674 & 0.0062 & 0.285         & 0.409      & 40.9×    \\
Isotonic Calibration        & 0.3439 & 0.0054 & 0.281         & 0.403      & 40.4×    
\end{longtblr}

\begin{multicols}{2}
Both Platt scaling and isotonic regression markedly improved probability
reliability without reducing ranking accuracy.

Although AUPRC remained nearly constant, the Brier score decreased by
nearly {\bfseries 17×}, confirming a major enhancement in the model's
probabilistic consistency.\\
Between the two techniques, isotonic regression achieved marginally
better local calibration (lower Brier) but incurred a slight global
ranking penalty (AUPRC ↓ 0.024).

Figure 1 shows that CatBoost outperforms all other models in AUPRC.
Figure 2 compares Brier scores, confirming superior calibration, while
Figure 3 presents reliability curves indicating near-linear alignment
between predicted and empirical probabilities, suggesting improved
probability interpretability.
\end{multicols}

\begin{figs}
    \fig[0.49\textwidth]{i2/image2}[Fig.1 - AUPRC by model on the IBM-AML dataset]
    \fig[0.49\textwidth]{i2/image3}[Fig.2 - Brier Score by model]
\end{figs}
\fig[0.45\textwidth]{i2/image4}[Fig.3 - Combined reliability diagram for all models]

To analyse the impact of temporal and dispersion-based behavioural
features, an ablation study was performed by sequentially excluding
specific feature groups. Table 3 summarises the results.

\tcap{Table 3 - Ablation results for behavioural features}
\begin{longtblr}[
  label = none,
  entry = none,
]{
  width = \linewidth,
  colspec = {Q[227]Q[90]Q[96]Q[529]},
  cells = {c},
  cells = {font = \small},
  row{1} = {font=\bfseries\small},
  cell{4}{2} = {font=\bfseries\small},
  cell{4}{3} = {font=\bfseries\small},
  hlines,
  vlines,
}
Removed Feature Group         & ΔAUPRC & ΔLift@1\% & Interpretation                                                                    \\
1-Day Rolling Window          & −0.012 & −2.0×     & Short-term volatility features improve responsiveness to rapid behaviour changes. \\
MAD (Mean Absolute Deviation) & −0.008 & −1.5×     & Dispersion-based features enhance anomaly sensitivity in low-signal regions.      \\
Both (1-Day + MAD)            & −0.023 & −4.0×     & Combined removal leads to noticeable degradation in ranking precision.            
\end{longtblr}

\begin{multicols}{2}
These results indicate that short-term (1-day) and dispersion (MAD)
features make a substantial contribution to AML detection sensitivity,
particularly in rare-event contexts.

Their removal caused a measurable degradation in both AUPRC and Lift@1
\%, indicating that behavioural temporal dynamics play an important role
in robust detection.

A backtesting experiment was conducted on the September 2022 period to
evaluate model robustness across time. To assess the temporal robustness
of the proposed models, a backtesting analysis was conducted on
the September 2022 transaction period. The results are presented
in Table 4 and visualised in Figure 4.
\end{multicols}

\tcap{Table 4 - Monthly backtesting results (AUPRC) on the IBM-AML dataset}

%% \begin{longtable}[]{@{}
%%   >{\centering\arraybackslash}p{(\linewidth - 6\tabcolsep) * \real{0.1576}}
%%   >{\centering\arraybackslash}p{(\linewidth - 6\tabcolsep) * \real{0.2761}}
%%   >{\centering\arraybackslash}p{(\linewidth - 6\tabcolsep) * \real{0.2740}}
%%   >{\centering\arraybackslash}p{(\linewidth - 6\tabcolsep) * \real{0.2924}}@{}}
%% \toprule\noalign{}
%% \begin{minipage}[b]{\linewidth}\centering
%% {\bfseries Month}
%% \end{minipage} & \begin{minipage}[b]{\linewidth}\centering
%% {\bfseries AUPRC (XGB)}
%% \end{minipage} & \begin{minipage}[b]{\linewidth}\centering
%% {\bfseries AUPRC (CAT)}
%% \end{minipage} & \begin{minipage}[b]{\linewidth}\centering
%% {\bfseries n (Transactions)}
%% \end{minipage} \\
%% \midrule\noalign{}
%% \endhead
%% \bottomrule\noalign{}
%% \endlastfoot
%% 2022-09 & 0.1784 & 0.3674 & 107 484 \\
%% \end{longtable}

\fig{i2/image5}[Fig. 4 - Temporal backtesting results
(AUPRC): comparison of CatBoost and XGBoost models on the September 2022
transaction period]

\begin{multicols}{2}
CatBoost maintained stable AUPRC ≈ 0.37 across time, outperforming
XGBoost (0.18). This temporal robustness indicates resilience to
distribution shifts, an important consideration for operational AML
settings.

However, analysis revealed clusters of false positives among accounts
exhibiting short-term transaction bursts (e.g., payroll disbursements or
high-frequency business activity). These legitimate but transient
patterns can be misclassified as laundering. Future work will
incorporate temporal decay factors and entity-type embeddings to
distinguish transient anomalies from systematic laundering behaviour.

The combined analysis supports the hypothesis that calibrated,
behaviour-aware models outperform both uncalibrated and purely
graph-based approaches. Behavioural rolling features enhance sensitivity
to dynamic patterns, while probability calibration ensures reliability
and interpretability of risk scores.

From an operational perspective, the calibrated CatBoost model
concentrates a substantially higher proportion of true positives within
the top-ranked transactions, which may support more efficient
prioritisation in AML review processes.

At the same time, these results should be interpreted as indicative, as
operational deployment would require additional validation, integration,
and regulatory assessment.

Nonetheless, the use of a synthetic dataset (IBM-AML) constitutes a
limitation, as it may not fully reflect the diversity and noise of
real-world financial data. Future research will expand this framework
using multi-institutional datasets, temporal graph neural networks
(TGNNs), and explainability methods (SHAP, LIME) to enhance
interpretability and cross-domain generalisation.

In summary, the results indicate that integrating temporal behavioural
modelling with probabilistic calibration improves ranking stability and
probability reliability in transaction-level AML detection. Behavioural
rolling features enhance sensitivity to dynamic patterns, while post-hoc
calibration improves the consistency of predicted risk scores.

At the same time, the reliance on a synthetic dataset (IBM-AML)
constitutes a limitation, and the reported findings should be
interpreted as indicative rather than directly transferable to
production environments. Further validation on real-world banking data,
along with the incorporation of temporal graph-based models and
explainability techniques, is necessary to assess practical
applicability.

A key limitation of this study is the reliance on a synthetic dataset
(IBM--AML), which, although widely used in academic research, may not
fully capture the heterogeneity, noise, and adaptive behaviour observed
in real-world financial systems.

As a result, the reported performance metrics should be interpreted as
indicative rather than directly transferable to production AML
environments.

Future work will therefore focus on validation using real or quasi-real
multi-institutional datasets.

Ethical considerations. This study uses a publicly available synthetic
dataset and does not involve personal or confidential financial records.
Nevertheless, automated AML decision support may introduce risks related
to false positives, biased prioritisation, and accountability;
therefore, future work on real-world deployment will include fairness,
explainability, and governance-oriented evaluation.

{\bfseries Conclusion.} This study presented a reproducible
machine-learning framework for transaction-level anti-money-laundering
(AML) detection that integrates temporal behavioural feature modelling
with probabilistic calibration. Experimental evaluation on the IBM
Anti-Money Laundering (IBM-AML) dataset demonstrated that calibrated
gradient boosting models, particularly CatBoost, achieve strong ranking
performance while improving probability, reliability, as reflected by a
low Brier score and stable operational metrics.

The results provide empirical evidence that the combination of rolling
behavioural features and post-hoc calibration enhances the robustness
and interpretability of AML risk scoring without degrading
discriminative performance. Short-term and dispersion-based behavioural
features were shown to be especially important in highly imbalanced
detection settings, while Platt scaling and isotonic regression
effectively aligned predicted scores with observed outcome frequencies.

From a methodological perspective, the study contributes an integrative
evaluation framework that jointly considers ranking quality,
calibration, and operational effectiveness, which are often analysed
separately in prior AML research. From a practical standpoint, the
findings highlight the potential value of calibration-aware and
behaviour-aware models for risk prioritisation in financial monitoring
systems.

At the same time, the reliance on a synthetic dataset constitutes a
limitation, and the reported results should be interpreted as indicative
rather than directly transferable to production environments. Future
work will focus on validation using real-world banking data,
incorporation of temporal graph-based models, and the application of
explainability techniques to further support transparent and auditable
AML decision-making.
\end{multicols}

\begin{center}
{\bfseries References}
\end{center}

\begin{refs}
1. Goldbarsht D., Benson K. From later to sooner: exploring compliance
with the global regime of anti-money laundering and counter-terrorist
financing in the legal profession //~Journal of Financial Crime.
-2023.-Vol.31(4). - P.795-809. DOI 10.1108/jfc-08-2023-0201.

2. Utama P.I., Novianto W.T., Purwadi H. The role of notaries/land deed
officials in combating money laundering crimes with the know your
customer principle //~International Journal of Educational Research \&
Social Sciences\emph{. -}2024.-Vol.5(2). - P.279-283. DOI
10.51601/ijersc. v5i2.801.

3. Reite E., Karlsen J., Westgaard E. Improving client risk
classification with machine learning to increase anti-money laundering
detection efficiency //~Journal of Money Laundering Control\emph{.}
-2025.-Vol.28(1). P.93-107.DOI 10.1108/JMLC-03-2024-0040.

4. Ali A., Abd Razak S., Othman S.H., Eisa T.A.E., Al-Dhaqm A., Nasser
M., Elhassan T., Elshafie H., Saif A. Financial fraud detection based
on machine learning: a systematic literature review //~Applied
Sciences\emph{.} 2022.- Vol.12(19). - Р.1-24. DOI 10.3390/app12199637.

5. Juneja S., Saurav A., Prabhakar A., Anand V., Goyal R., Thakur V.
Enhancing fraud detection with gradient boosting algorithms in credit
card transactions//Proc. of the 2025 Int. Conf. on Cognitive Computing
in Engineering, Communications, Sciences and Biomedical Health
Informatics (IC3ECSBHI). -2025.-P.1-8. DOI
10.1109/IC3ECSBHI63591.2025.10990658.

6. Lopata A., Gudas S., Butleris R., Rudžionis V., Žioba L., Veitaitė I.,
Dilijonas D., Grišius E., Zwitserloot M. Financial data anomaly
discovery using behavioral change indicators //~Electronics\emph{. -}
2022. Vol.11(10). -Р.1-14. DOI 10.3390/electronics11101598.

7. Weber M., Domeniconi G., Chen J., Weidele D.K.I., Bellei C., Robinson
T., Leiserson C.E. Anti-money laundering in bitcoin: experimenting
with graph convolutional networks for financial forensics//arXiv
preprint\emph{.-}2019.-Р.1-7.DOI 10.48550/arXiv.1908.02591.

8. Johannessen F., Jullum M. Finding money launderers using heterogeneous
graph neural networks//arXiv preprint\emph{.} -2023.-Р.1-20. DOI:
10.48550/arXiv.2307.13499.

9. Wan F., Li P. A novel money laundering prediction model based on a
dynamic graph convolutional neural network and long short-term
memory//Symmetry\emph{.} -2024.-Vol.16.-Р.1-22. DOI
10.3390/sym16030378.

10. Poon C.-H., Kwok J., Chow C., Choi J.-H. LineMVGNN: anti-money
laundering with line-graph-assisted multi-view graph neural networks
//AI\emph{.} -2025.-Vol.6(4). DOI 10.3390/ai6040069.

11. IBM/AML-Data. IBM Transactions for Anti-Money Laundering (AML)
Dataset {[}Online
resource{]}.-2019.URL:\href{https://www.kaggle.com/datasets/ealtman2019/ibm-transactions-for-anti-money-laundering-aml.-\%20Date}{https://www.kaggle.com/datasets/ealtman2019/ibm-transactions-for-anti-money-laundering-aml.-
Date} of access: 10.10.2025.

12. Adelino J.M., Tierra M.C.M., Villafuerte M.Q. Development and
evaluation of multi-week tropical cyclone strike probability forecasts
in the Philippines //~Philippine Journal of Science\emph{.}
-2023.-Vol.152(SI). -P.1-13.DOI 10.56899/152.s1.09.

13. Xu P., Davoine F., Zha H., Denœux T. Evidential calibration of binary
SVM classifiers //Approximate Reasoning. -2016.-Vol.-72.-P.55-70. DOI
10.1016/j.ijar.2015.05.002.

14. Nguyen H.D., Fryer D. A binary-response regression model based on
support vector machines //~arXiv preprint\emph{.} -2020. P.1-20. DOI
10.48550/arXiv.2007.11902.

15. Silva Filho T., Song H., Perello-Nieto M., Santos-Rodriguez R., Kull
M., Flach P. Classifier calibration: a survey on how to assess and
improve predicted class probabilities//Machine Learning\emph{.
-}2023.-Vol.112. P.1-50. DOI 10.1007/s10994-023-06336-7.

16. Chen T., Guestrin C. XGBoost: A scalable tree boosting system
//~Proceedings of the 22nd ACM SIGKDD International Conference on
Knowledge Discovery and Data Mining (KDD 2016)\emph{.} - New York: ACM,
2016. - P.785-794. - DOI 10.1145/2939672.2939785.

17. Sahin E. K. Assessing the predictive capability of ensemble tree
methods for landslide susceptibility mapping using XGBoost, gradient
boosting machine, and random forest //~SN Applied Sciences\emph{.} -
2020. - Vol.2:1308. DOI 10.1007/s42452-020-3060-1.

18. Prokhorenkova L., Gusev G., Vorobev A., Dorogush A. V., Gulin A.
CatBoost: Unbiased boosting with categorical features //~Advances in
Neural Information Processing Systems (NeurIPS 2018). - 2018. -
P.6638--6648. DOI
\href{https://doi.org/10.48550/arXiv.1706.09516}{10.48550/arXiv.1706.09516}

19. Shwartz-Ziv R., Armon A. Tabular data: Deep learning is not all you
need //~Information Fusion\emph{. -} 2022. - Vol.81. - P.84-90. DOI
10.1016/j.inffus.2021.11.011.

20. Weber M., Domeniconi G., Chen J., Weidele D. K. I., Bellei C.,
Robinson T., Leiserson C. E. Anti-money laundering in Bitcoin:
Experimenting with graph convolutional networks for financial forensics
//~arXiv preprint\emph{.} - 2019.- arXiv:1908.02591. DOI
10.48550/arXiv.1908.02591.

21. Johannessen F., Jullum M. Finding money launderers using
heterogeneous graph neural networks // arXiv preprint. - 2023. -
arXiv:2307.13499. DOI 10.48550/arXiv.2307.13499.

22. Hamilton W., Ying Z., Leskovec J. Inductive representation learning
on large graphs//arXiv preprint\emph{. -}2017.- P.1-19. DOI
10.48550/arXiv.1706.02216.

23. Probst P., Wright M. N., Boulesteix A.-L. Hyperparameters and tuning
strategies for random forest //~Wiley Interdisciplinary Reviews: Data
Mining and Knowledge Discovery. - 2019.- Vol.9(3): e1301. DOI
10.1002/widm.1301.

24. Nyberg O., Klami A. Reliably calibrated isotonic regression
//~Advances in Knowledge Discovery and Data Mining: PAKDD 2021\emph{.}
Lecture Notes in Computer Science. - Cham: Springer. -2021.- Vol.12712.
- P.578--589. DOI 10.1007/978-3-030-75762-5\_46.

25. Böken B. On the appropriateness of Platt scaling in
classifier calibration //~Information Systems\emph{.} - 2021. - Vol.
95:101641. DOI
\href{https://doi.org/10.1016/j.is.2020.101641}{10.1016/j.is.2020.101641}
\end{refs}

\begin{info}
\emph{{\bfseries Information about the authors}}

Azbergen T.K. - master's student, Kazakh-British Technical University,
Almaty, Kazakhstan, e-mail: azbergen.torebek@mail.ru;

Tleubayeva A. O. - doctoral student, Astana IT University, Astana,
Kazakhstan, e-mail: Arailym.tll@gmail.com;

Aituov A. T. - PhD, Kazakh-British Technical University, Almaty,
Kazakhstan, e-mail: a.aituov@kbtu.kz;

Sabyrzhan Y. T. - master of Technical Sciences, Astana IT University,
Astana, Kazakhstan, e-mail: sabyrzhan.yermakhan@astanait.edu.kz.

\emph{{\bfseries Сведения об авторах}}

Азберген Т. К. - магистрант, Казахстанско-Британский технический
университет, Алматы, Казахстан, e-mail: azbergen.torebek@mail.ru;

Тлеубаева А. О.- докторант, Astana IT University, Астана, Казахстан,
e-mail: Arailym.tll@gmail.com;

Айтуов А. Т. - PhD, Казахстанско-Британский технический университет,
Алматы, Казахстан, e-mail: a.aituov@kbtu.kz;

Сабыржан Е. Т. - магистр технических наук, Astana IT University, Астана,
Казахстан, e-mail: sabyrzhan.yermakhan@astanait.edu.kz.
\end{info}
